% Front matter of Volume One: edition notice (PUB-REF-0001 AN-20, fixed
% order; boilerplate newly written, VO-12) and the volume Introduction with
% Conventions Used in This Manual (AN-21 F-ATE, NT-35; annex Sec 7, 8.1).
% The motto follows the edition notice (operator directive 2026-10-17).
\Gates{test-pubs-build}
\begin{TitleVerso}
{\imLegalFont
\textbf{The following paragraph does not apply to any country where such
provisions are inconsistent with local law:} INITECH SYSTEMS CORPORATION
PROVIDES THIS PUBLICATION ``AS IS'' WITHOUT WARRANTY OF ANY KIND, EITHER
EXPRESS OR IMPLIED, INCLUDING, BUT NOT LIMITED TO, THE IMPLIED WARRANTIES OF
MERCHANTABILITY AND FITNESS FOR A PARTICULAR PURPOSE.  Some states do not
allow the exclusion of implied warranties, so the above exclusion may not
apply to you.\par
This publication may contain technical inaccuracies or typographical errors.
The information in it is revised from time to time, and the revisions are
published in new editions.  Initech Systems Corporation may change the
products and programs described in this publication at any time without
notice.\par
Requests for copies of Initech publications and for technical information
about Initech products should be made to your authorized Initech dealer or
your Initech representative.\par}

{\footnotesize
Reader's Comment Forms are provided at the back of this publication.  If the
forms have been removed, address comments to: \InitechAddressLine.
\draft{PUB-REF-0001 O-5 and D-12: the address and permit number are to be
checked against real businesses before print}
Initech may use or distribute the information you supply in any way it
believes appropriate without incurring any obligation to you.\par
The software described in this publication is furnished under the Initech
Program License Agreement and may be used and copied only as that agreement
permits.  The sample databases, report forms and programs on the Sample
Programs and Utilities diskette may be copied and changed for use with
InitechBase.\par}
\IsThisGoodForTheCompany
\end{TitleVerso}

\FrontChapter{Introduction}
InitechBase\tm{} is a database management system for the InitechOS\tm{}
personal computer.  It keeps your business records in database files, finds
and lists them, changes them, totals them in reports, and runs the programs
you write to do the same work automatically.  You give InitechBase its
instructions as commands typed at the dot prompt, or as programs that it
reads from a file.

\section{The InitechBase Documentation}
The documentation is divided into sections.  Each section has its own
letter, which begins every page number in the section: page L3-7 is the
seventh page of Chapter~3 of \doctitle{Learning InitechBase}, and UI-2 is the
second page of the introduction to \doctitle{Using InitechBase}.

\SideHead{Getting Started (G)}Describes the InitechBase package, tells you
how to copy InitechBase to your fixed disk, and shows you how to start it and
how to leave it.  Read it first.

\SideHead{Learning InitechBase (L)}A tutorial.  In nine chapters you use the
sample databases of the Initech staff roster to list, find, change, total and
report on records, and you write your first program.  You need no experience
of databases.

\SideHead{Using InitechBase (U)}The reference.  It describes the parts of a
command, the rules for expressions, and the InitechBase environment; then it
gives every command, every function and every error message in alphabetical
order.  Use it when you need the exact form of a command.

\SideHead{Programming with InitechBase (P)}In Volume Two.  It teaches you to
write InitechBase programs by building the Initech Accounting System, the
sample accounts receivable application, a module at a time.

\SideHead{Index (X)}The index at the back of each volume covers both
volumes.  Each page reference begins with its section letter.

\SideHead{Quick Reference}The \doctitle{InitechBase Quick Reference Guide}
lists the syntax of every command and function on a few pages that fit beside
your keyboard.

\section{Conventions Used in This Manual}
\SideHead{What you type}Text that you type is shown in bold type:
\type{USE STAFF}.  You may type commands in capital or small letters;
InitechBase does not distinguish between them.  The manual prints commands in
capitals.

\SideHead{Keys}Keys are named as they are printed on the keys of your
keyboard: \key{Enter}, \key{Backspace}, \key{Esc}.  Two keys joined by a
hyphen are pressed together: hold down \key{Ctrl} and press \key{C} for
\key{Ctrl-C}.  At the end of every command line you press \key{Enter}; the
steps of the tutorial say so.

\SideHead{Examples}The examples are printed on a shaded panel exactly as they
appear on your screen.  Every line you type at the dot prompt follows the
prompt, a period and a space, and is printed in bold:
\begin{TypeIn}
. USE STAFF
\end{TypeIn}
Every example in this manual was run against the files of the Sample
Programs and Utilities diskette, and the panel shows what InitechBase
displayed.

\SideHead{Syntax}The form of each command is given as a syntax line.  Words
in capitals are typed as shown.  Angle brackets enclose an item you supply,
such as \xp{expC}, a character expression; you do not type the brackets.
Square brackets enclose an item you may leave out.  A slash separates items
of which you choose one.  The symbols are listed in full at the beginning of
Chapter~5 of \doctitle{Using InitechBase}.

\SideHead{Boxes}A TIP box gives a quicker or easier way to do something.  A
NOTE box gives information you need in a particular case.  A WARNING box
tells you of a step that can destroy data.

\SideHead{Programs}Program listings are printed in bold on a shaded panel,
as they are stored in the program files of the sample diskette.
